\section{Literature Review}
\label{sec:literature_review}

\subsection{Semantic SLAM for UAVs}
\label{subsec:semantic_slam}

Semantic SLAM for aerial platforms must be lightweight, robust to GPS signal loss, and capable of producing useful semantic maps for planning and multi-robot coordination. Liu et al. proposed SlideSLAM -- a sparse, lightweight, decentralized metric-semantic SLAM for multi-robot navigation that uses an object-based representation and enables efficient map merging in heterogeneous aerial-ground teams~\cite{liu2024slideslam}. Tao et al. integrated active exploration with metric-semantic SLAM for resource-constrained aerial robots, explicitly balancing exploration and localization uncertainty (including semantic loop closures), which informs S4D-TAM strategies for uncertainty-aware mapping in GNSS-denied environments~\cite{tao2024active}. Liu et al. presented a real-time semantic SLAM that detects and models tree trunks and planes from LiDAR to constrain poses and enable drift compensation during long-range autonomous flights under forest canopy, directly relevant to drift mitigation and semantic constraints in S4D-TAM trajectory planning~\cite{liu2022largescale}.

Tian et al. proposed SAME -- a ground-air collaboration system that fuses semantic Octree units from a UAV and a UGV into a coherent global map supporting collaborative exploration, offering a semantic octree architecture and fusion approach useful in S4D-TAM for cross-platform mapping and multi-perspective scene reconstruction~\cite{tian2024same}. Xiang and Su developed a LiDAR-based semantic mapping pipeline that creates hierarchical semantic maps together with occupancy grids for navigation and task adaptation, illustrating practical semantic mapping designs and hierarchical semantics that S4D-TAM can exploit for task-oriented scene understanding~\cite{xiang2021customizable}.

Li et al. introduced Hi-SLAM -- a hierarchical categorical representation within 3D Gaussian Splatting for compact encoding of multiple semantic classes and improved semantic mapping and tracking, directly applicable to S4D-TAM requirements for scalable semantic scene representations~\cite{li2024hislam, keetha2024splatam}. More recent semantic SLAM works employing 3D Gaussian Splatting include SGS-SLAM~\cite{hu2024sgsslam}, GS4/GaussianSLAM~\cite{yugay2023gaussianslam}, SemGauss-SLAM~\cite{zou2024semgaussslam}, Photo-SLAM~\cite{huang2024photoslam}, DynaGS-SLAM~\cite{wu2024dynags}, GSFF-SLAM~\cite{pan2024gsffslam}, OVO-SLAM~\cite{haghighi2024ovoslam}, MSGS-SLAM~\cite{li2024msgsslam}, NIS-SLAM~\cite{li2024nisslam}, and NEDS-SLAM~\cite{dai2024nedsslam}, demonstrating growing interest in real-time dense semantic representations for SLAM.

\subsection{Point Cloud Transformers}
\label{subsec:point_cloud_transformers}

Point cloud transformer architectures and related models provide expressive spatiotemporal encoding of 3D geometry and dynamics, useful for the point-cloud-based perception and temporal prediction modules of S4D-TAM. PCPNet introduces a transformer network for future LiDAR point cloud prediction with a semantic auxiliary training objective for semantically consistent forecasts, directly applicable to S4D-TAM for point-level prediction and semantic alignment~\cite{pcpnet2024}. Silva et al. proposed S2TPVFormer, which extends TPV embeddings with hybrid cross-view temporal attention for temporally consistent semantic 3D occupancy, offering a spatiotemporal transformer design that S4D-TAM can adopt for temporally consistent semantic occupancy mapping~\cite{silva2024s2tpvformer}.

Chen et al. presented AdaOcc, which employs adaptive forward-view transformations and flow modeling (with strong Swin camera backbone models) to improve voxel occupancy and flow prediction, providing practical camera-to-3D transformer strategies applicable in vision-oriented S4D-TAM modules~\cite{chen2024adaocc}. Work on point cloud forecasting as a proxy for 4D occupancy prediction redefines the problem toward sensor-agnostic spatiotemporal (4D) occupancy forecasting, directing S4D-TAM toward world-oriented 4D occupancy targets for robust temporal modeling~\cite{yu2025driveoccworld}.

Further relevant works include Point-BERT~\cite{yu2022pointbert}, which introduces masked point modeling for pre-training point cloud transformers; SAT3D / Point Attention Network~\cite{ji2020point} for semantic segmentation of 3D scenes using slot attention; and YOGO~\cite{zhang2021you} and Efficient Point Transformer~\cite{he2024efficient, guo2021pct}, demonstrating efficient point cloud processing with token representation and relation inference modules relevant to S4D-TAM token compression.

\subsection{4D Occupancy Forecasting}
\label{subsec:occupancy_forecasting}

Accurate 4D occupancy forecasting provides the spatiotemporal scene model required by S4D-TAM for risk assessment and trajectory planning. Yu et al. proposed Drive-OccWorld, which adapts a vision-centric 4D world model with memory normalization and action-conditioned decoding for occupancy and flow prediction, and integrates forecasting with end-to-end planning, providing an approach that S4D-TAM can leverage to couple occupancy prediction with trajectory selection~\cite{yu2025driveoccworld}. Zheng et al. presented OccWorld -- a spatiotemporal generative transformer architecture for 4D occupancy prediction, demonstrating effectiveness on driving datasets and offering S4D-TAM a transformer backbone for long-horizon volumetric prediction~\cite{zheng2024occworld}.

Ma et al. provided Cam4DOcc -- a standardized camera-only benchmark and baselines for current and future occupancy estimation, providing S4D-TAM with evaluation protocols and camera-only baseline designs for 4D occupancy forecasting~\cite{ma2023cam4docc}. OccSora uses a 4D scene tokenizer and a diffusion transformer for generative long-horizon 4D occupancy simulation conditioned on trajectories, offering S4D-TAM a generative world simulator approach for scenario generation and planning evaluation~\cite{wang2024occsora}. Liao et al. proposed I2-World with hierarchical intra/inter tokenizers and an encoder-decoder design for efficient compression and generation of temporally consistent 4D occupancy, presenting a tokenization strategy that S4D-TAM can apply for compact world representations and fast inference~\cite{liao2025i2world}.

Zheng et al. presented GaussianAD, which represents scenes through semantic 3D Gaussians refined from images, predicts 3D flows for those Gaussians, and integrates forecasting with planning -- this sparse semantic Gaussian approach is directly relevant to S4D-TAM's lightweight scene abstractions and prediction-based planning~\cite{zheng2024gaussianad}.

\subsection{World Models and Planning}
\label{subsec:world_models}

World models unify perception, prediction, and planning; for S4D-TAM, they provide a mechanism for simulating the future, evaluating trajectories, and incorporating action-conditioned forecasts. Chen et al. proposed DriveWorld -- a Memory State-Space Model with dynamic memory and static scene propagation for spatiotemporal pre-training, demonstrating downstream gains in detection, mapping, tracking, forecasting, and planning -- relevant to S4D-TAM pre-training strategies and multi-task world modeling~\cite{chen2024driveworld}.

Zhao et al. used a contrastive world model objective to learn visual representations for drone navigation and visuomotor policy learning, illustrating how world model pre-training can improve navigation policies, which S4D-TAM can leverage for improved visual generalization~\cite{zhao2023learning}. Mei et al. proposed IR-WM, which focuses on predicting residual changes relative to a temporal BEV prior and couples forecasting with planning, offering S4D-TAM a compact residual prediction approach that reduces redundant static background modeling and improves planning integration~\cite{mei2025irwm}.

Cai et al. presented NEUSIS -- a compositional neuro-symbolic framework combining neuro-symbolic perception, a probabilistic world model with Bayesian filtering, and hierarchical search planning, exemplifying an interpretable and uncertainty-aware world modeling and planning workflow useful for S4D-TAM mission-level decision-making~\cite{cai2024neusis}. Liu et al. proposed HPP, which integrates occupancy-conditioned prediction with game-theoretic motion forecasting and planning, demonstrating methods for aligning prediction patterns with planning objectives that S4D-TAM can adopt for improved trajectory selection under uncertainty and interaction~\cite{liu2025hybrid}.

Feng et al. provided a comprehensive survey of world models for autonomous driving, covering 4D occupancy forecasting, generative simulation, and planning integration, providing a conceptual foundation and taxonomy that can guide S4D-TAM design choices in world model selection, pre-training, and evaluation~\cite{feng2025survey}.

\subsection{UAV Navigation in GNSS-Denied Environments}
\label{subsec:gnss_denied}

Autonomous UAV navigation in GNSS-degraded environments requires multi-sensor fusion, robust SLAM algorithms, and uncertainty-aware planning strategies. Yue et al. proposed an air-ground cooperative UAV autonomous navigation system for GPS-denied environments that combines visual odometry, IMU, and communication with a ground robot for robust localization~\cite{yue2025airground}. Radwan et al. presented UAV-assisted visual SLAM generating reconstructed 3D scene graphs in GPS-denied environments, demonstrating scene graph representations for scene understanding and planning~\cite{radwan2024uav}.

Nieuwenhuisen et al. developed an autonomous navigation system for micro aerial vehicles in complex GNSS-denied environments, using LiDAR, IMU, and obstacle-aware planning~\cite{nieuwenhuisen2015autonomous, nieuwenhuisen2016autonomous}. Yue et al. proposed semantically guided visual autonomous navigation for UAVs that employs semantic segmentation to identify navigable areas and obstacles for improved path planning~\cite{yue2024semantic}. Hu et al. provided a review of vision-based UAV localization and optimization-based path planning in GPS-denied environments, synthesizing SLAM, object detection, and trajectory planning methods~\cite{hu2025review}.

Further relevant works include UAV navigation systems using multi-sensor fusion (IMU, camera, LiDAR, radar)~\cite{delmerico2018benchmark, qin2018vinsmono, shan2020liosam, zhang2014loam, xu2022fastlio2, faessler2015automatic}, KAZE feature-based SLAM~\cite{alcantarilla2013kaze}, map-aided navigation and dead reckoning~\cite{brossard2020aiimu}, and platforms for drone racing and agile navigation~\cite{delmerico2019uzhfpv}. Castano presented autonomous UAV navigation in GPS-denied environments using LiDAR point clouds for localization and mapping~\cite{castano2023autonomous}.
